\section{Introduction}
\label{sec:intro}

Retrieval-augmented generation (RAG) systems assume that the retriever surfaces evidence that is
\emph{valid for the question being asked}. For temporally grounded questions---``Why was there a
jump in transactions per block in 2023?''---validity depends on time: a passage about the right
topic but the wrong period is misleading evidence. The RETECO shared task
\citep{abdallah2026tempo,ali2026recor} operationalises this with TEMPO, a 13-domain benchmark of
StackExchange questions whose relevant passages must be both topical and temporally aligned.

The official BM25 baseline reaches only 0.108--0.118 nDCG@10 on Track~1, far below its
0.44 on the conversational Track~2, confirming that ``temporal grounding is not a keyword
problem''. We hypothesise three reasons: (1)~queries often express time \emph{implicitly}
(``after the halving'', ``two weeks ago''), (2)~relevant passages frequently share few terms
with the question, and (3)~nothing in a bag-of-words ranking prefers documents from the
right period.

\sys addresses each: \todo{one sentence per component, mapped to (1)--(3)}.

Our contributions are:
\begin{itemize}
  \item A metadata layer that attaches normalised temporal profiles, entities and short
        descriptions to 1.17M documents and exposes them to every retrieval and reranking stage.
  \item \emph{Guidance distillation}: using TEMPO's train-split query guidance as demonstrations so
        an LLM emits the benchmark's own temporal annotation schema for unseen queries.
  \item A six-channel hybrid retriever (sparse, dense, HyDE, sub-query, temporal knowledge graph)
        with weighted RRF and a temporal-alignment prior, plus a temporal-aware listwise reranker.
  \item Ablations and per-domain / per-intent analyses on all 13 domains.
\end{itemize}
